\scanpage{380}
\mspage{119}
% archive sheet number: 379
si la loi de composition admet l'objet final de $T$, supposé exister, comme unité bilatère). Dans ce cas, $\underline t=(t_n)$ et $\underline t'=(t'_n)$ sont reliés entre eux par $\underline t*\underline t'=(t_n*t'_n)_n$, \deleted{\unsure{et}} \inserted{puisqu'}on a un morphisme de systèmes projectifs
\diagram{\begin{tikzcd}[ampersand replacement=\&,row sep=normal,column sep=small]
\&\underline t*\underline t'\arrow[dl]\arrow[dr]\&\\
\underline t\&\&\underline t'\,.
\end{tikzcd}}

Je ne sais d'ailleurs pas si la flèche dans $\mathrm{Hot}$
\[
C_\infty(\underline t)\mathrel{\unsure{-|\to}}C_\infty(\underline t')
\]
associée à $\underline t'\mathrel{-}\underline t$ dans $\underline{\mathrm{Hom\,cof}}$ \deleted{\unsure{de}} ne dépend que de $\underline t_{\unsure{0|\circ}}$, $\underline t'$, et pas du choix de $\underline t'\longrightarrow\underline t$. Ce qui est clair, \deleted{\unsure{c'est que dans le cas}} \inserted{c'est que son image} dans $\deleted{\unsure{T=\underline{\mathrm{Tg}}|T=\underline{T_g}}}$ $\mathrm{Hot}_0^{\unsure{\wedge|1}}$ est unique, c'est l'identité quand on identifie $C_\infty(\underline t)$ et $C_\infty(\underline t')$ au seul et $=$ élément de $\mathrm{Hot}_0^\wedge$ défini par $\underline{\lim}\,H_T$. \deleted{On doit avoir}

Et pourtant, je m'aperçois à l'instant que, sous l'hypothèse de l'existence de la loi de composition et de son unité bilatère, et admettant de plus que la loi en question est associative à isom. près, \uline{on a transitivité} dans les iso. (canoniques !) entre les $\underline C_\infty(t)$ qu'on vient d'indiquer !
